\documentclass[11pt]{article}
\usepackage{amsmath,amssymb,amsthm,hyperref}
\newtheorem{theorem}{Theorem}
\newtheorem{lemma}[theorem]{Lemma}
\title{Exchange Refinements of Boolean Reciprocal Features}
\author{}
\date{November 23, 2026}
\begin{document}
\maketitle

\paragraph{Scope.}
This note derives a merged, two-step averaging identity and proves
rigidity of its edge weights. The subset-inclusion lemmas are kernel
checked in \texttt{Results/SubsetInclusionRigidity.lean}; the reciprocal
identities and polynomial degree arguments are manual proofs. This note
supplies no universal head upper bound.

Let $b_0,\ldots,b_n>0$, let $z_0=1$, and let $z_i=x_i\in\{0,1\}$
for $i>0$. Set
\[
 T=\sum_{i=0}^n b_i,\qquad B(x)=\sum_{i=0}^n b_i z_i,
 \qquad K_b(x)=\frac{b_0}{B(x)}.
\]
For $0<|\delta|<\min_i b_i$ and $i<j$ define
\[
 E_{ij}(x)=K_{b+\delta(e_i-e_j)}(x)
          +K_{b-\delta(e_i-e_j)}(x)-2K_b(x).
 \tag{1}
\]
Every weight vector in (1) is strictly positive. The terms are
admissible constant-numerator heads, and repeated occurrences of
$K_b$ can be merged into one head.

\begin{theorem}[Exact exchange averaging]
For every Boolean input,
\[
                  \sum_{0\le i<j\le n} b_i b_j E_{ij}(x)=0.
 \tag{2}
\]
\end{theorem}
\begin{proof}
The homogeneous version of the P\'olya identity is
\[
 K_b=\sum_{i=0}^n\frac{b_i}{T}K_{b+\delta e_i}.
 \tag{3}
\]
It follows by grouping selected and unselected colors, or by the
one-step calculation in \texttt{polya-reciprocal-audit.tex}.
The same identity holds for a negative increment whenever all
resulting weights remain positive. Apply (3) with increment $\delta$
and then at each child with increment $-\delta$. The probability
of the ordered pair $(i,j)$ is
\[
 \frac{b_i(b_j+\delta\mathbf1_{i=j})}{T(T+\delta)}.
\]
For $i=j$ the endpoint is $b$. For $i\ne j$ it is
$b+\delta(e_i-e_j)$. Subtract $K_b$ from the resulting average,
multiply by $T(T+\delta)$, and group the two orders of each
unordered pair. This gives (2). All denominators and transition
probabilities used in these two steps are positive.
\end{proof}

There are at most $n(n+1)+1$ distinct denominator forms in (2).
The count does not assert that these forms are all different.

\begin{theorem}[Rigidity of common-step edge weights]
Suppose $n\ge1$ and let $w_{ij}\in\mathbb R$ for
$0\le i<j\le n$. Under the same assumptions on $b$ and $\delta$,
\[
 \sum_{i<j} w_{ij}E_{ij}(x)=0\quad\hbox{for every Boolean }x
 \tag{4}
\]
if and only if there is a real $\lambda$ such that
$w_{ij}=\lambda b_i b_j$ for every edge. In particular, signed
weights do not add freedom to this common-step identity.
\end{theorem}
\begin{proof}
Direct subtraction gives, for $1\le i<j\le n$,
\[
 E_{ij}=\frac{2\delta^2 b_0(x_i-x_j)^2}
                    {B(B^2-\delta^2)},\qquad
 E_{0j}=\frac{2\delta^2(b_0-B)(1-x_j)}
                    {B(B^2-\delta^2)}.
 \tag{5}
\]
For example, when $x_i=x_j$ the first expression is zero; when
they differ its two shifted denominators are $B+\delta$ and
$B-\delta$. In the second expression, the case $x_j=1$ is zero,
whereas the case $x_j=0$ also shifts the numerator by $\pm\delta$.
Because $B\ge b_0>|\delta|$, the common denominator in (5) is
strictly positive. Thus (4) is equivalent to vanishing on the cube
of the multilinear quadratic polynomial
\[
 P(x)=b_0\sum_{1\le i<j\le n}w_{ij}(x_i+x_j-2x_ix_j)
      -\left(\sum_i b_i x_i\right)\sum_j w_{0j}(1-x_j).
 \tag{6}
\]
Here and below diagonal squares are reduced using $x_i^2=x_i$.
Its coefficient of $x_ix_j$, for distinct nonzero indices, is
\[
                  -2b_0 w_{ij}+b_iw_{0j}+b_jw_{0i}.
 \tag{7}
\]
Its coefficient of $x_i$ is
\[
                  b_0\sum_{j\ne i}w_{ij}
                         -b_i\sum_{j\ne i}w_{0j}.
 \tag{8}
\]
Vanishing of a multilinear quadratic on the cube forces these
coefficients to be zero: first evaluate at singleton supports and
then at two-element supports. Set $s=\sum_{j=1}^n b_j>0$ and
$u=\sum_{j=1}^n w_{0j}$. Substitute (7) into twice (8). The
result is
\[
                         s w_{0i}-b_i u=0.
 \tag{9}
\]
Hence $w_{0i}=b_i u/s$, and (7) gives
$w_{ij}=b_i b_j u/(b_0s)$. Taking $\lambda=u/(b_0s)$ proves
necessity, including $n=1$, where there are no nonanchor edges.
Sufficiency follows from (2).
\end{proof}

\paragraph{Consequences for the attempted construction.}
The two urn steps merge their return paths and produce an exact
Boolean cancellation, but choosing an arbitrary graph of weights
does not generate a larger family of cancellations. A nonzero
common-step relation uses every edge, with precisely the weights
in (2), up to one scalar. The next result also controls a bounded
number of increment magnitudes. Nonconstant numerator weights and
general merged refinement graphs remain outside these statements.

\begin{lemma}[Two elementary inclusion matrices]
Let $1\le k\le n-1$. If real $a_i$ satisfy
$\sum_{i\in S}a_i=0$ for every $k$-element subset $S$ of $[n]$,
then every $a_i$ is zero. If $2\le k\le n-2$ and real symmetric
edge weights $a_{ij}$ satisfy
$\sum_{\{i,j\}\subseteq S}a_{ij}=0$ for every $k$-element
subset $S$, then every $a_{ij}$ is zero.
\end{lemma}
\begin{proof}
For the first assertion, compare subsets with the same $k-1$
elements and with respective last elements $i$ and $j$. This makes
all $a_i$ equal; a $k$-element sum then makes their common value
zero. For the second, fix distinct $i,j$ and compare the sums for
$R\cup\{i\}$ and $R\cup\{j\}$, where $R$ ranges over all
$(k-1)$-element subsets outside $\{i,j\}$. The first assertion
on those $n-2$ indices gives $a_{ir}=a_{jr}$ for every remaining
$r$. Thus all incident edge weights at each vertex agree. Symmetry
then makes all edge weights equal. A $k$-element sum, containing
$\binom{k}{2}>0$ edges, makes their common value zero.
\end{proof}

The two assertions are kernel checked as
\texttt{HeadComplexity.fixed\_card\_subset\_sum\_zero} and
\texttt{HeadComplexity.fixed\_card\_pair\_sum\_zero}. The second
uses a symmetric double sum, which counts each unordered pair twice.

\begin{lemma}[Separated quadratic coefficients]
Let $B=b_0+\sum_i b_ix_i$ with every $b_i\ne0$, let $M\ge1$,
and assume $n\ge2M+2$. Let $P_0,\ldots,P_{M-1}$ be multilinear
polynomials of degree at most two with $P_j(0)=0$. If
\[
                \sum_{j=0}^{M-1} B(x)^{2j}P_j(x)=0
                       \quad(x\in\{0,1\}^n),
 \tag{10}
\]
then all $P_j$ are zero.
\end{lemma}
\begin{proof}
Use descending induction on the largest index $j$ still present,
and put $r=2j$. Write the quadratic coefficients of $P_j$ as
$p_{ab}$. In the multilinear reduction of $B^rP_j$, the
coefficient at a subset $S$ of size $r+2$ is
\[
 r!\left(\prod_{i\in S}b_i\right)
       \sum_{\{a,b\}\subseteq S}\frac{p_{ab}}{b_ab_b}.
 \tag{11}
\]
To reach that degree, all $r$ factors from $B$ must choose
distinct coordinates outside the two from $P_j$. This proves
(11), including its factor $r!$. Lower-index terms in (10)
have degree at most $r$, so these coefficients vanish. We have
$2\le r+2\le2M\le n-2$; the second inclusion assertion makes
every $p_{ab}$ zero.

Now let $p_a$ be the linear coefficients of $P_j$. At degree
$r+1$, the analogous coefficient is
\[
 r!\left(\prod_{i\in S}b_i\right)
                          \sum_{a\in S}\frac{p_a}{b_a}.
 \tag{12}
\]
Again the lower-index terms have degree at most $r$. Since
$1\le r+1\le n-1$, the first inclusion assertion makes every
$p_a$ zero. The constant term is zero by hypothesis. Remove
$P_j$ and repeat. Vanishing of all multilinear coefficients used
here follows from uniqueness of the multilinear representative
on the cube, or directly by induction on support size.
\end{proof}

\begin{theorem}[Few-magnitude exchange rigidity]
Assign each edge a nonzero increment $\delta_{ij}$ with
$|\delta_{ij}|<\min_i b_i$, and let $E_{ij}$ use that increment
in (1). Suppose their absolute values belong to a set of $M$
distinct positive numbers and $n\ge2M+2$. If
$\sum_{i<j}w_{ij}E_{ij}=0$ on the cube, then either all $w_{ij}$
are zero, or every edge has the same increment magnitude and
$w_{ij}=\lambda b_i b_j$ for one nonzero scalar $\lambda$.
\end{theorem}
\begin{proof}
List the distinct magnitudes as $d_1,\ldots,d_M$. For each
magnitude, form the polynomial $P_s$ in (6), setting edge weights
outside that class to zero. Formula (5) gives
\[
             \sum_{s=1}^M\frac{d_s^2 P_s(x)}{B(x)^2-d_s^2}=0.
 \tag{13}
\]
Clear the positive denominators. The resulting expression is a
linear combination of $1,B^2,\ldots,B^{2M-2}$ whose coefficients
are quadratic polynomials vanishing at zero. The preceding lemma
makes all of these coefficient polynomials zero. Thus, as a
polynomial identity in a separate variable $u$,
\[
       \sum_{s=1}^M d_s^2 P_s(x)
                         \prod_{t\ne s}(u-d_t^2)=0.
 \tag{14}
\]
Substitute $u=d_s^2$. The magnitudes are positive and distinct,
so $P_s=0$ for every $s$. The common-step rigidity theorem
applied to each class says that its edge weights are
$\lambda_s b_i b_j$ on the complete graph. If $\lambda_s\ne0$,
all these weights are nonzero, so that class must contain every
edge. Otherwise every class is zero. This proves the assertion.
\end{proof}

\paragraph{Small-dimensional exceptions and finite checks.}
The dimension restriction in the last theorem cannot simply be
discarded. For $n=3$, take $b=(3,3,3,3)$. Give the edges
$\{0,3\}$ and $\{1,2\}$ increment magnitude $2$ and weight
$77/320$; give the other four edges magnitude $1$ and weight $1$.
This is a positive nonconstant-magnitude relation on all eight
vertices. Indeed the quadratic polynomial (6) for the indicated
matching vanishes at Hamming weights zero, one, and three.
The complete-graph polynomial is zero, so its complement has
the opposite polynomial. On the remaining layer $B=9$, and
\[
             \frac{77}{320}\frac4{81-4}=\frac1{81-1}.
\]
This proves the claimed relation using (13).

The checker \texttt{check\_polya\_exchange.py} independently
evaluates (1) with rational arithmetic, verifies (5) and (2), and
checks common-step matrix ranks in dimensions one through seven.
It also searches the two magnitudes $1,2$ for two specified positive
weight vectors. A full-rank matrix modulo the stated prime proves
full rank over the rationals; every deficient candidate used as an
exception is recomputed and verified over the rationals. Exhaustion
for those fixed choices in dimensions four and five does not prove
rigidity for arbitrary weights or arbitrary increment magnitudes.

\paragraph{Separate literature check.}
Seshadri and Weso\l{}owski, \emph{Martingales Defined by Reciprocals
of Sums and Related Characterizations} (2004), Theorems 2.1 and 3.1,
classify certain inverse-sum regression identities for positive iid
random variables. Their gamma, inverse Gaussian, Kendall, and
Borel--Tanner identities concern conditional expectations. They do
not give a bound on the number of affine reciprocal atoms needed
to realize arbitrary Boolean signs. The exchange identity above
is proved directly and does not invoke those characterization results.
\url{https://pages.mini.pw.edu.pl/~wesolowskij/publikacje/2004SesWes.pdf}

\end{document}
